\section*{第 \ref{ch:b2:surfaces} 章——曲面}

\begin{solution}{exo:b2:surfaces:1}
令 $\sigma(\theta, r) = (r\cos\theta,\ r\sin\theta,\ r)$：
\[
\sigma_\theta = (-r\sin\theta,\ r\cos\theta,\ 0),
\qquad
\sigma_r = (\cos\theta,\ \sin\theta,\ 1),
\]
当 $r > 0$ 时两者线性无关。点 $M_0 =
\sigma(\theta_0, r_0)$ 处的\omterm{def:b2:surfaces:tangent}{切平面}过 $M_0$，方向为
$\sigma_\theta, \sigma_r$。而 $M_0 - O = r_0(\cos\theta_0,
\sin\theta_0, 1) = r_0\,\sigma_r(\theta_0, r_0)$ 本身就是切方向，
故顶点 $O$ 位于\omterm{def:b2:surfaces:tangent}{切平面}上。（这
是锥面的一般性质：锥面由过顶点的直线
构成，且\omterm{def:b2:surfaces:tangent}{切平面}包含经过
切点的直母线。）
\end{solution}

\begin{solution}{exo:b2:surfaces:2}
将\cref{prop:b2:surfaces:gradient}用于 $f(x,y,z) =
\frac{x^2}{a^2} + \frac{y^2}{b^2} + \frac{z^2}{c^2}$，则 $\nabla
f(x_0,y_0,z_0) = 2\bigl(\frac{x_0}{a^2}, \frac{y_0}{b^2},
\frac{z_0}{c^2}\bigr) \neq 0$（于曲面上）。\omterm{def:b2:surfaces:tangent}{切平面}为
\[
\frac{x_0}{a^2}(x - x_0) + \frac{y_0}{b^2}(y - y_0)
+ \frac{z_0}{c^2}(z - z_0) = 0,
\qquad\text{即}\qquad
\frac{x_0\,x}{a^2} + \frac{y_0\,y}{b^2} + \frac{z_0\,z}{c^2} = 1,
\]
利用 $(x_0, y_0, z_0)$ 满足椭球面方程——
“拆分平方”法则，推广了球面\omterm{def:b2:surfaces:tangent}{切平面}方程 $\langle
M_0, M \rangle = R^2$。
\end{solution}

\begin{solution}{exo:b2:surfaces:3}
$\sigma_u = (-v\sin u,\ v\cos u,\ a)$ 与 $\sigma_v = (\cos u,\
\sin u,\ 0)$，故
\[
E = v^2 + a^2, \qquad F = 0, \qquad G = 1,
\qquad
\sqrt{EG - F^2} = \sqrt{v^2 + a^2} > 0
\]
（螺旋面处处正则，包括其轴线上 $v =
0$）。这一片的\omterm{def:b2:surfaces:area}{面积}为：
\[
\mathcal{A} = \int_0^{2\pi}\!\!\int_0^1
\sqrt{v^2 + a^2}\;\dd v\,\dd u
= 2\pi\cdot\frac12\Bigl[v\sqrt{v^2 + a^2}
+ a^2\ln\bigl(v + \sqrt{v^2 + a^2}\bigr)\Bigr]_0^1 ,
\]
即 $\mathcal{A} = \pi\Bigl(\sqrt{1 + a^2} +
a^2\ln\frac{1 + \sqrt{1 + a^2}}{a}\Bigr)$。
\end{solution}

\begin{solution}{exo:b2:surfaces:4}
偏导数为：
\[
\sigma_\theta = \bigl(-(R + r\cos\psi)\sin\theta,\
(R + r\cos\psi)\cos\theta,\ 0\bigr),
\qquad
\sigma_\psi = \bigl(-r\sin\psi\cos\theta,\
-r\sin\psi\sin\theta,\ r\cos\psi\bigr).
\]
于是
\[
E = (R + r\cos\psi)^2,
\qquad
F = r\sin\psi\cos\psi\,(R + r\cos\psi)
\bigl(\sin\theta\cos\theta - \sin\theta\cos\theta\bigr) = 0,
\qquad
G = r^2 ,
\]
故 $\sqrt{EG - F^2} = r(R + r\cos\psi) \geq r(R - r) > 0$，
处处正则。\omterm{def:b2:surfaces:area}{面积}为：
\[
\mathcal{A}
= \int_0^{2\pi}\!\!\int_0^{2\pi} r(R + r\cos\psi)
\,\dd\theta\,\dd\psi
= 2\pi r \Bigl(2\pi R + r\int_0^{2\pi}\cos\psi\,\dd\psi\Bigr)
= 4\pi^2 R r ,
\]
其中 $\cos\psi$ 项的积分为零——帕普斯定理：\omterm{def:b2:surfaces:area}{面积}
$=$（旋转圆周的长）$\times$（
其中心所经过的距离）。
\end{solution}

\begin{solution}{exo:b2:surfaces:5}
对 $\sigma(x, y) = (x, y, f(x,y))$，有 $\sigma_x = (1, 0, f_x)$，
$\sigma_y = (0, 1, f_y)$，故 $E = 1 + f_x^2$，$F = f_xf_y$，$G = 1
+ f_y^2$，并有
\[
EG - F^2 = (1 + f_x^2)(1 + f_y^2) - f_x^2f_y^2
= 1 + f_x^2 + f_y^2 ,
\]
由此得所述\omterm{def:b2:surfaces:area}{面积}公式。对 $f = \frac12(x^2 + y^2)$，在
单位圆盘上 $1 + f_x^2 + f_y^2 = 1 + x^2 + y^2$，并在极
坐标下（$x = \rho\cos\alpha$，$y = \rho\sin\alpha$，雅可比\omterm{def:b2:linalg:det}{行列式}
$\rho$，见\cref{ch:b2:multint}）：
\[
\mathcal{A}
= \int_0^{2\pi}\!\!\int_0^1 \sqrt{1 + \rho^2}\;\rho
\,\dd\rho\,\dd\alpha
= 2\pi\Bigl[\tfrac13(1 + \rho^2)^{3/2}\Bigr]_0^1
= \frac{2\pi}{3}\bigl(2\sqrt2 - 1\bigr).
\]
\end{solution}

\begin{solution}{exo:b2:surfaces:6}
由\cref{ex:b2:surfaces:spherearea}的计算，$E =
R^2\cos^2\varphi$，$F = 0$，$G = R^2$，故由
\cref{prop:b2:surfaces:curvelength}可得
\[
L = \int_a^b R\sqrt{\cos^2\varphi(t)\,\theta'(t)^2
+ \varphi'(t)^2}\;\dd t .
\]
设端点为 $(\theta_0, \varphi_1)$ 与 $(\theta_0,
\varphi_2)$，$\varphi_1 < \varphi_2$。对任一连接两点的曲线，
\[
L \geq \int_a^b R\,\abs{\varphi'(t)}\,\dd t
\geq R\,\Bigl|\int_a^b \varphi'(t)\,\dd t\Bigr|
= R\,(\varphi_2 - \varphi_1),
\]
略去非负项 $\cos^2\varphi\,\theta'^2$，并利用
积分的三角不等式。子午线弧
$\theta \equiv \theta_0$，$\varphi$ 从 $\varphi_1$ 增至
$\varphi_2$ 的弧，其长度恰为 $R(\varphi_2 - \varphi_1)$，它
是最短的。（子午线是大圆；这是最初、
最基本的情形：球面测地线皆为
大圆。）
\end{solution}

\begin{solution}{exo:b2:surfaces:7}
设 $\gamma$ 为 $S$ 上画出的一条曲线，令 $g(t) = \norm{\gamma(t)
- \Omega}^2$，则 $g'(t) = 2\langle \gamma'(t),\ \gamma(t) -
\Omega\rangle$。过 $M = \gamma(t)$ 的法线经过
$\Omega$（由假设），故 $\gamma(t) - \Omega$ 是一个\emph{法向量}，
与\omterm{def:b2:surfaces:tangent}{切平面}正交，特别地与
速度 $\gamma'(t)$（见\cref{prop:b2:surfaces:velocity}）：$g' =
0$，从而 $g$ 沿每条画出的曲线为常数。

考虑集合 $S_c = \{M \in S : \norm{M - \Omega}^2 = c\}$，
它是 $S$ 中的闭集；它也是 $S$ 中的\omterm{def:b2:metric:topology}{开集}：在其任一点附近，$S$
是一个正则函数图像，故 $S$ 中任一邻近点都可
用一条画出的曲线（一段提升的线段）与之相连，沿该曲线 $g$ 为常数。由于
$S$ 连通，$S_c$ 在 $c$ 取适当值时非空，故 $S = S_c
\subseteq$ 以 $\Omega$ 为球心、$\sqrt c$ 为半径的球面
（见\cref{ch:b2:metric}：连通性论证）。
\end{solution}

\begin{solution}{exo:b2:surfaces:8}
记 $\Phi(u', v') = (u, v)$。由链式法则，
\[
\tilde\sigma_{u'} = \frac{\partial u}{\partial u'}\sigma_u
+ \frac{\partial v}{\partial u'}\sigma_v,
\qquad
\tilde\sigma_{v'} = \frac{\partial u}{\partial v'}\sigma_u
+ \frac{\partial v}{\partial v'}\sigma_v ,
\]
其中 $\sigma$ 的偏导数均在 $\Phi(u',v')$ 处取值。将
向量积双线性展开，并利用 $\sigma_u \wedge \sigma_u =
\sigma_v \wedge \sigma_v = 0$，$\sigma_v \wedge \sigma_u =
-\sigma_u \wedge \sigma_v$：
\[
\tilde\sigma_{u'} \wedge \tilde\sigma_{v'}
= \Bigl(\frac{\partial u}{\partial u'}
\frac{\partial v}{\partial v'}
- \frac{\partial v}{\partial u'}
\frac{\partial u}{\partial v'}\Bigr)\,
\sigma_u \wedge \sigma_v
= \det J_\Phi \cdot (\sigma_u \wedge \sigma_v)\circ\Phi .
\]
取模即得恒等式。再由
\cref{ch:b2:multint}中的变量替换公式，对映射 $\Phi$ 在 $K' =
\Phi^{-1}(K)$ 上应用，得：
\[
\iint_{K'} \norm{\tilde\sigma_{u'} \wedge \tilde\sigma_{v'}}
\,\dd u'\dd v'
= \iint_{K'} \norm{\sigma_u \wedge \sigma_v}\circ\Phi\;
\abs{\det J_\Phi}\,\dd u'\dd v'
= \iint_K \norm{\sigma_u \wedge \sigma_v}\,\dd u\,\dd v :
\]
两种参数化对同一块
曲面片给出相同的\omterm{def:b2:surfaces:area}{面积}。
\end{solution}

\begin{solution}{exo:b2:surfaces:9}
在球坐标卡中，$z = R\sin\varphi$，球带
对应于 $\varphi_1 \leq \varphi \leq \varphi_2$，其中
$z_i = R\sin\varphi_i$。取\omterm{def:b2:surfaces:area}{面积}元
$R^2\cos\varphi\,\dd\theta\,\dd\varphi$
（见\cref{ex:b2:surfaces:spherearea}）：
\[
\mathcal A = \int_{\varphi_1}^{\varphi_2}\!\!\int_0^{2\pi}
R^2\cos\varphi\,\dd\theta\,\dd\varphi
= 2\pi R^2(\sin\varphi_2 - \sin\varphi_1)
= 2\pi R\,(z_2 - z_1) .
\]
结果仅依赖于高度 $z_2 - z_1$：等厚的
切片\omterm{def:b2:surfaces:area}{面积}相等，无论切在赤道
或切在极点处——阿基米德帽盒定理，也正是
外切圆柱的侧\omterm{def:b2:surfaces:area}{面积}（$2\pi R \cdot
2R = 4\pi R^2$）与球面\omterm{def:b2:surfaces:area}{面积}相等的原因。
\end{solution}

\begin{solution}{exo:b2:surfaces:10}
$\sigma_\theta = (-r\sin\theta,\ r\cos\theta,\ 0)$，且
$\sigma_z = (r'\cos\theta,\ r'\sin\theta,\ 1)$，于是
\[
\sigma_\theta \wedge \sigma_z
= (r\cos\theta,\ r\sin\theta,\ -r\,r') ,
\]
这是点 $M = (r\cos\theta, r\sin\theta, z)$ 处的法向量。过该点的
法线为
\[
t \mapsto \bigl(r(1 + t)\cos\theta,\ r(1 + t)\sin\theta,\
z - t\,r\,r'\bigr),
\]
当 $t = -1$ 时到达 $(0,\ 0,\ z + r(z)\,r'(z))$：每条
法线都与轴相交，交点高度为 $z + rr'$。（这
是三维空间中旋转对称性在法向场里仍然保持的
原因。）
\end{solution}

\begin{solution}{exo:b2:surfaces:11}
$\sigma_u = (-\sin u, \cos u, 0)$，$\sigma_v = (0, 0, 1)$：$E
= 1$，$F = 0$，$G = 1$。由
\cref{prop:b2:surfaces:curvelength}，画出曲线的长度为
$\int\sqrt{u'^2 + v'^2}\,\dd t$——即其参数影子
$(u(t), v(t))$ 在平面中的长度：该坐标卡是
局部等距映射（圆柱的展开）。螺线
$t \mapsto (\cos t, \sin t, ct)$，$t \in \intcc0{2\pi}$，其
影子为从 $(0,0)$ 到 $(2\pi, 2\pi c)$ 的线段，长度为
$2\pi\sqrt{1 + c^2}$。任一条由 $(1,0,0)$ 到 $(1,
0, 2\pi c)$ 绕一整圈的画出曲线，其影子\omterm{def:b2:metric:continuity}{连续}
地连接 $(0, 0)$ 与 $(2\pi, 2\pi c)$，平面长度 $\geq$
直线段的长度；由于长度相等，螺线
最短。
\end{solution}

\begin{solution}{exo:b2:surfaces:12}
函数 $g(M) = \norm{M}^2$ 在紧的
$S$ 上\omterm{def:b2:metric:continuity}{连续}，故在某点 $M_0$ 取得最大值。
对每条曲线 $\gamma$（画在 $S$ 上，且 $\gamma(0) = M_0$），函数 $t
\mapsto \norm{\gamma(t)}^2$ 在 $t = 0$ 处取极大值，故其
导数 $2\langle\gamma'(0), M_0\rangle$ 为零，故 $M_0$
与每个切向量正交，亦即 $S$ 在
$M_0$ 处的法向。由于 $\nabla f(M_0) \neq 0$ 也给出法线方向
（见\cref{prop:b2:surfaces:gradient}），$\nabla f(M_0)$ 与
$\vect{OM_0}$ 共线。对椭球面，
径向性意味着
\[
\Bigl(\frac{x_0}{a^2}, \frac{y_0}{b^2},
\frac{z_0}{c^2}\Bigr) = \mu\,(x_0, y_0, z_0) :
\]
各坐标满足 $x_0(\frac1{a^2} - \mu) = 0$ 等；
由于 $a^{-2}, b^{-2}, c^{-2}$ 互不相同，至多
一个坐标非零，且曲面上的解为
六个轴端点 $(\pm a, 0, 0)$、$(0, \pm b, 0)$、$(0,
0, \pm c)$。最远的点是 $(\pm a, 0, 0)$，它们
距原点的距离为 $a = \max(a,b,c)$。
\end{solution}

\begin{solution}{pb:b2:surfaces:1}
\textbf{1.} 展开，并利用 $A$ 的对称性
（$t^{\mathsf T}\!APY = (At)^{\mathsf T}PY$）：
\[
q(PY + t) = Y^{\mathsf T}P^{\mathsf T}\!APY +
2\,(At + b)^{\mathsf T}PY + \bigl(t^{\mathsf T}\!At +
2b^{\mathsf T}t + c\bigr),
\]
这正是上式所示的三项。$A' = P^{\mathsf T}\!AP =
P^{-1}AP$ 与 $A$ 相似，\omterm{def:b2:reduction:charpoly}{特征多项式}相同，
谱、秩、\omterm{thm:b2:quadratic:sylvester}{符号差}相同。最后，$\{q = 0\} = \{sq = 0\}$
对 $s \neq 0$ 成立，故集合只与方程\emph{在相差一个数乘的意义下}
相联系；以 $s$ 放缩时，所有\omterm{def:b2:reduction:eigen}{特征值}都乘以
$s$。

\textbf{2.} 由谱定理，存在 $P \in O(3)$，使
$P^{\mathsf T}\!AP = \operatorname{diag}(\lambda_1, \lambda_2,
\lambda_3)$；在第 1 小题中取 $t = 0$，方程化为
$\sum\lambda_iy_i^2 + 2\sum\beta_iy_i + c = 0$，其中 $\beta =
P^{\mathsf T}b$。

\textbf{3.} 当 $\lambda_i \neq 0$ 时，$\lambda_iy_i^2 +
2\beta_iy_i = \lambda_i\bigl(y_i +
\frac{\beta_i}{\lambda_i}\bigr)^2 -
\frac{\beta_i^2}{\lambda_i}$；作平移 $z_i = y_i +
\beta_i/\lambda_i$（$z_i = y_i$ 对 $i > r$ 成立），得
\[
\sum_{i \leq r}\lambda_iz_i^2 + 2\sum_{i > r}\beta_iz_i +
c'' = 0, \qquad
c'' = c - \sum_{i\leq r}\frac{\beta_i^2}{\lambda_i}.
\]

\textbf{4.} $q(2\Omega - X) = (2\Omega - X)^{\mathsf T}
A(2\Omega - X) + 2b^{\mathsf T}(2\Omega - X) + c$；展开
并减去 $q(X)$，二次项相消，且
\[
q(2\Omega - X) - q(X) = -4\,\langle A\Omega + b,\ X\rangle +
4\,\langle A\Omega + b,\ \Omega\rangle .
\]
若 $A\Omega + b = 0$，则上式恒为零：此时点
反射保持 $q$，从而保持该\omterm{pb:b2:surfaces:1}{二次曲面}。反之，
“$q(2\Omega - X) = q(X)$ 对一切 $X$ 成立”说明上面的仿射
函数在整个 $\R^3$ 上恒为零，这迫使它的
线性部分 $A\Omega + b$ 为零。于是中心 $=$ 下式的解集
$A\Omega = -b$：其解集非空当且仅当 $b \in \operatorname{im}A$
（以 $\ker A$ 为方向的仿射子空间），而恰为一点
当且仅当 $A$ 可逆。

\textbf{5.} \omterm{thm:b2:quadratic:sylvester}{符号差} $(3,0)$（所有 $\lambda_i > 0$）：当 $\delta >
0$ 时给出 $\frac{x^2}{a^2} + \frac{y^2}{b^2} + \frac{z^2}{c^2}
= 1$，其中 $a = \sqrt{\delta/\lambda_1}$ 等——一个
椭球面；$\delta = 0$：单点 $O$；$\delta < 0$：
空集。\omterm{thm:b2:quadratic:sylvester}{符号差} $(2,1)$（$\lambda_1, \lambda_2 > 0 >
\lambda_3$）：当 $\delta > 0$ 时，$\frac{x^2}{a^2} + \frac{y^2}{b^2}
- \frac{z^2}{c^2} = 1$，为单叶双曲面；$\delta =
0$ 时为锥面 $\frac{x^2}{a^2} + \frac{y^2}{b^2} =
\frac{z^2}{c^2}$；当 $\delta < 0$ 时，$\frac{z^2}{c^2} -
\frac{x^2}{a^2} - \frac{y^2}{b^2} = 1$，为双叶
双曲面（$\abs z \geq c$ 时分为两支）。

\textbf{6.} 二次部分的矩阵 $A$ 对角元为
$1$，非对角元为 $2$：$A = 2J - I$。由于 $J$ 的谱为
$\{3, 0, 0\}$（$(1,1,1)$ 为对应于 $3$ 的\omterm{def:b2:reduction:eigen}{特征向量}），$A$ 的
谱为 $\{5, -1, -1\}$，\omterm{def:b2:reduction:eigen}{特征值} $5$ 对应的特征空间为
$\R(1,1,1)$。在旋转后的坐标下：$5u^2 - v^2 - w^2 =
1$，即 $\frac{u^2}{1/5} - v^2 - w^2 = 1$：双叶
双曲面，由于\omterm{def:b2:reduction:eigen}{特征值} $-1$ 相等，它是绕
轴 $\R(1,1,1)$ 的旋转双曲面。

\textbf{7.} 该曲面为 $\bigl(\frac xa - \frac
zc\bigr)\bigl(\frac xa + \frac zc\bigr) = \bigl(1 - \frac
yb\bigr)\bigl(1 + \frac yb\bigr)$。对 $(\lambda : \mu) \neq
(0:0)$ 定义直线
\[
D_{\lambda:\mu} :\quad
\lambda\Bigl(\frac xa - \frac zc\Bigr) = \mu\Bigl(1 - \frac
yb\Bigr), \qquad
\mu\Bigl(\frac xa + \frac zc\Bigr) = \lambda\Bigl(1 + \frac
yb\Bigr)
\]
（两个独立的仿射方程，表示一条直线）。将
两式相乘可知，$D_{\lambda:\mu}$ 上每一点都在
曲面上（当 $\lambda\mu \neq 0$ 时）；$\lambda = 0$
或 $\mu = 0$ 的情形可直接验证（例如 $\lambda = 0$ 时：$y =
b$，$\frac xa = -\frac zc$，满足原方程）。
第二族 $D'_{\lambda:\mu}$ 交换两个右端
的因子。

\textbf{8.} 固定曲面上一点 $M$。$M \in
D_{\lambda:\mu}$ 的条件构成一个 $2\times2$ 齐次线性方程组，
关于 $(\lambda, \mu)$ 的，其\omterm{def:b2:linalg:det}{行列式}为
\[
\Bigl(\frac{x^2}{a^2} - \frac{z^2}{c^2}\Bigr) - \Bigl(1 -
\frac{y^2}{b^2}\Bigr) = 0
\]
恰因 $M$ 位于\omterm{pb:b2:surfaces:1}{二次曲面}上，故存在非平凡
解 $(\lambda : \mu)$。系数矩阵
不为零（否则将迫使 $1 - \frac yb = 1 + \frac yb =
0$），因此其秩为 $1$，解在相差一个数乘的
意义下唯一：族中恰有一条直线经过 $M$。
对第二族同理，且两条直线
互异（在 $(a, 0, 0)$ 处分别为 $\{x = a,\ \frac zc = \frac
yb\}$ 与 $\{x = a,\ \frac zc = -\frac yb\}$）：单叶
双曲面是双重直纹面。

\textbf{9.} 设 $M = (x, y, z)$ 为双曲面上一点，$\rho^2 =
\frac{x^2}{a^2} + \frac{y^2}{b^2} = 1 + \frac{z^2}{c^2} \geq
1$。取点 $N = (x,\ y,\ \varepsilon c\rho)$，其中
$\varepsilon$ 为 $z$ 的符号，且满足 $\frac{x^2}{a^2} +
\frac{y^2}{b^2} - \frac{(c\rho)^2}{c^2} = 0$，故 $N \in C$，且
\[
d(M, C) \leq \abs{z - \varepsilon c\rho}
= c\bigl(\rho - \sqrt{\rho^2 - 1}\bigr)
= \frac{c}{\rho + \sqrt{\rho^2 - 1}} .
\]
若 $\norm M \to \infty$，则 $\rho \to \infty$（三个
坐标都被 $\rho$ 的倍数控制），故 $d(M, C)
\to 0$。截面 $x = a$：$\frac{y^2}{b^2} -
\frac{z^2}{c^2} = 0$，为
第 8 小题中的一对相交直线——在 $x = -a$ 处同理。

\textbf{10.} 取 $\lambda_3 = 0$，由第 3 题得到
$\lambda_1z_1^2 + \lambda_2z_2^2 + 2\beta_3z_3 + c'' = 0$。在
特征基中，$\operatorname{im}A = \operatorname{Vect}(e_1,
e_2)$，由第 4 题，中心存在当且仅当 $\beta_3 = 0$。若
$\beta_3 \neq 0$ 时：作平移 $z_3 \mapsto z_3 -
c''/(2\beta_3)$ 消去常数项，剩下 $\lambda_1z_1^2
+ \lambda_2z_2^2 + 2\beta_3z_3 = 0$，即经重命名后为 $z_3 = px^2 +
qy^2$，为\emph{椭圆抛物面}，当
$\lambda_1\lambda_2 > 0$ 时；为\emph{双曲抛物面}，当
$\lambda_1\lambda_2 < 0$ 时——显然无中心。若
$\beta_3 = 0$ 时：方程 $\lambda_1z_1^2 + \lambda_2z_2^2 +
c'' = 0$ 不含 $z_3$，此\omterm{pb:b2:surfaces:1}{二次曲面}为相应平面二次曲线的柱面：
椭圆柱面、直线，或空集（$\lambda_1\lambda_2 > 0$ 时）；
双曲柱面，
或两相交平面（$\lambda_1\lambda_2 < 0$ 时）——
并且中心构成整条直线
$\{(z_1^*, z_2^*)\} \times \R$。

\textbf{11.} $xy - z = 0$。代入 $x = \frac{u +
v}{\sqrt2}$，$y = \frac{u - v}{\sqrt2}$（旋转角为 $\pi/4$）：
$xy = \frac{u^2 - v^2}2$，于是方程化为 $z =
\frac12(u^2 - v^2)$，这是双曲抛物面，
即图中的马鞍面（相差因子 $\frac12$）。

\textbf{12.} 直线 $\{x = x_0,\ z = x_0y\}$（以 $y$ 为参数）
显然位于 $z = xy$ 上，同样 $\{y = y_0,\ z =
xy_0\}$ 也在其上；过点 $(x_0, y_0, x_0y_0)$ 的恰是这两条。
唯一性：若 $t \mapsto (x_0 + tv_1, y_0 + tv_2, z_0 +
tv_3)$ 仍在曲面上，则 $t^2$ 的系数（取自 $(x_0
+ tv_1)(y_0 + tv_2) - z_0 - tv_3$）给出 $v_1v_2 = 0$，于是
$v_1 = 0$ 或 $v_2 = 0$，故直线落在两族之一中：
每点上两族各有一条直线经过，
即第二种双重直纹\omterm{pb:b2:surfaces:1}{二次曲面}。

\textbf{13.} 配方：$(x - 1)^2 + (y + 2)^2 = 2$，
对 $z$ 不加限制：半径为
$\sqrt2$ 的正圆柱面，轴为竖直直线 $\{(1, -2, z) : z \in
\R\}$ ——正如第 10 题所预言，这是一条中心直线。

\textbf{14.} 取 $\lambda_2 = \lambda_3 = 0$，约化方程为
$\lambda_1z_1^2 + 2\beta_2z_2 + 2\beta_3z_3 + c''
= 0$。对核平面 $(z_2, z_3)$ 作旋转，使线性形式得以对齐：
$2\beta_2z_2 + 2\beta_3z_3 = 2\beta w$，其中
$\beta = \sqrt{\beta_2^2 + \beta_3^2}$。若 $\beta \neq 0$，
则平移 $w$ 以吸收 $c''$：$\lambda_1z_1^2 + 2\beta w =
0$，为\emph{抛物柱面}；若 $\beta = 0$：
$\lambda_1z_1^2 = -c''$ 给出两平行平面（$c''
\lambda_1 < 0$），二重平面（$c'' = 0$），或
空集。对 $(x + y)^2 = z$：令 $u = \frac{x + y}{\sqrt2}$，方程化为
$z = 2u^2$：抛物柱面，
在沿 $(1, -1, 0)$ 方向的平移下不变。

\textbf{15.} 每个\omterm{pb:b2:surfaces:1}{二次曲面}经旋转与平移，必化为下列之一：
（秩 3）椭球面、点、空集、
单叶双曲面、锥面、双叶双曲面；
（秩 2）椭圆抛物面、双曲抛物面、椭圆
柱面、直线、空集、双曲柱面、两相交平面；
（秩 1）抛物柱面、两平行平面、
二重平面、空集。将
三种空集的变体视为不同的\emph{仿射类型}，
则方程的类型共计十七种；其中九种为真正的曲面：
椭球面、两种双曲面、锥面、两种
抛物面，以及三种柱面。

\textbf{16.} 算法。(i) 读出 $(A, b, c)$；计算
$A$ 的\omterm{def:b2:reduction:charpoly}{特征多项式}，其\omterm{def:b2:reduction:eigen}{特征值}（由谱定理，皆为实数），
并求 $r = \operatorname{rank}A$。(ii)
解 $A\Omega = -b$（高斯消元法）：是否可解，
即是否存在中心。(iii) 若可解，则平移至中心：
方程化为 $\sum\lambda_iz_i^2 + c'' = 0$，其中 $c''
= c + \langle b, \Omega\rangle$；按 $r$、\omterm{thm:b2:quadratic:sylvester}{符号差}
及 $c''$ 的符号，利用第 5、10、14 题分类。(iv) 若不可解，
即 $r \leq 2$，则仿照第 10、14 题旋转并化简，
得抛物面（$r = 2$）或抛物柱面（$r = 1$），
椭圆型或双曲型，取决于
$\lambda_1\lambda_2$ 的符号。每一步都是有限次计算：
实系数三次方程的根、秩、线性方程组。

\textbf{17.} $A$ 的对角元为 $1$，非对角元为 $-1$：$A = 2I
- J$，谱为 $\{2 - 3,\ 2,\ 2\} = \{-1, 2, 2\}$，其中 $-1$
的\omterm{def:b2:reduction:eigen}{特征子空间}为 $\R(1,1,1)$。\omterm{thm:b2:quadratic:sylvester}{符号差}为 $(2,1)$，$b = 0$，右端项
$\delta = 1 > 0$：$2u^2 + 2v^2 - w^2 = 1$，为
单叶双曲面，即绕轴 $\R(1,1,1)$ 的旋转双曲面。

\textbf{18.} 配方：$(x-1)^2 + (y+2)^2 - (z-1)^2 +
(-1 - 4 + 1 + 4) = 0$，即
\[
(x-1)^2 + (y+2)^2 = (z-1)^2 :
\]
常数项消去——这是一个正圆\emph{锥面}，
顶点（亦即唯一中心）为 $(1, -2, 1)$，轴平行于
$Oz$。

\textbf{19.} 二次部分 $x^2 + 4xy + y^2$ 的矩阵为
$\begin{psmallmatrix}1 & 2\\ 2 & 1\end{psmallmatrix}$（位于
$xy$ 平面内），\omterm{def:b2:reduction:eigen}{特征值}为 $3$（属于 $(1,1)$）及 $-1$（属于
$(1,-1)$）。令 $u = \frac{x+y}{\sqrt2}$，$v =
\frac{x-y}{\sqrt2}$，则其等于 $3u^2 - v^2$，于是曲面
方程为
\[
z = \tfrac32u^2 - \tfrac12v^2 :
\]
这是双曲抛物面（$A$ 的秩为 $2$，$b$ 具有
沿 $\ker A = \R e_z$ 的分量：无中心）。

\textbf{20.} 刚体运动将方程的数据按
$A \mapsto P^{\mathsf T}\!AP$ 变换（\omterm{def:b2:reduction:eigen}{特征值}不变），而
标准形除了坐标置换之外，没有其他残余自由度，
以及将整个方程乘以一个标量（取
$> 0$ 以保持形式不变）。两个中心标准形
在等距意义下重合，当且仅当系数列表在
置换及公共正因子意义下一致；对椭球面而言，
即当半轴 $(a, b, c)$ 相同时。仿射意义下，亦可分别对各坐标作伸缩
（$z_i \mapsto z_i\sqrt{\abs
{\lambda_i}}$），这样\omterm{def:b2:reduction:eigen}{特征值}被消去，只剩下其符号：
每个椭球面都化为 $u^2 + v^2 + w^2 = 1$，
即球面。西尔维斯特惯性定理
（\cref{thm:b2:quadratic:sylvester}）保证\omterm{thm:b2:quadratic:sylvester}{符号差}
在任何可逆线性变换下保持不变：第 15 题中的仿射类型
确实互不相同。

\textbf{21.} 仿射地参数化平面：$M = P + su +
tv$。于是 $q(P + su + tv)$ 是次数 $\leq 2$ 的多项式
（将\omterm{def:b2:quadratic:def}{二次型}双线性展开），变量为 $(s, t)$，故
截面 $\{q = 0\}$ 是平面上的一条二次曲线，可能
退化。对 $z = xy$ 与平面 $z = c$：$xy = c$，
当 $c \neq 0$ 时为双曲线，而当 $c = 0$ 时为两条坐标
直线，即过原点的一对直母线。

\textbf{22.} \emph{椭球面}：有界，不含直线。
\emph{双叶双曲面} $\frac{z^2}{c^2} -
\frac{x^2}{a^2} - \frac{y^2}{b^2} = 1$：限制到 $p + tv$ 上；
$t^2$ 的系数 $\frac{v_3^2}{c^2} - \frac{v_1^2}{a^2}
- \frac{v_2^2}{b^2}$ 须为零，故 $v_3 \neq 0$（否则 $v =
0$）；$t$ 的系数给出 $\frac{p_3v_3}{c^2} =
\frac{p_1v_1}{a^2} + \frac{p_2v_2}{b^2}$，再由柯西–施瓦茨
不等式可得
\[
\frac{p_3^2}{c^2}
= \frac{c^2}{v_3^2}\Bigl(\frac{p_1v_1}{a^2} +
\frac{p_2v_2}{b^2}\Bigr)^2
\leq \frac{c^2}{v_3^2}\Bigl(\frac{p_1^2}{a^2} +
\frac{p_2^2}{b^2}\Bigr)\frac{v_3^2}{c^2}
= \frac{p_1^2}{a^2} + \frac{p_2^2}{b^2},
\]
于是常数项 $\leq 0 \neq 1$，故不含直线。
\emph{椭圆抛物面} $z = \frac{x^2}{a^2} +
\frac{y^2}{b^2}$：$t^2$ 的系数迫使 $v_1 = v_2 =
0$，此时方程关于 $t$ 是非常数的一次式，不含直线。
\emph{锥面} $x^2 + y^2 = z^2$：其上直线满足 $q(p) =
q(v) = B(p, v) = 0$（对洛伦兹型）；
平面柯西–施瓦茨不等式取等号：$\abs{p_1v_1 + p_2v_2} = \abs{p_3v_3} =
\sqrt{(p_1^2 + p_2^2)(v_1^2 + v_2^2)}$ 迫使 $(p_1, p_2)
\parallel (v_1, v_2)$，进而 $p = kv$：所有直线
都过顶点——只有一族。\emph{柱面}：对
椭圆柱面与抛物柱面，$t^2$ 的系数迫使
$v_1 = v_2 = 0$（仅有直母线）；对双曲柱面
$\frac{x^2}{a^2} - \frac{y^2}{b^2} = 1$，$\frac{v_1}a = \pm
\frac{v_2}b$，$v_2 \neq 0$，经 $t$ 的系数
推出 $\frac{p_1^2}{a^2} = \frac{p_2^2}{b^2}$，与
常数项 $1$ 矛盾：此时仍只有竖直直母线。于是
双重直纹的\omterm{pb:b2:surfaces:1}{二次曲面}恰为单叶双曲面
与双曲抛物面。

\textbf{23.} 用高斯配方法：$x^2 - 2xy - 2zx = (x - y - z)^2 - y^2 -
z^2 - 2yz$，于是
\[
q = (x - y - z)^2 - 4yz = (x - y - z)^2 + (y - z)^2 - (y +
z)^2 :
\]
三个独立平方，符号为 $(+, +, -)$——\omterm{thm:b2:quadratic:sylvester}{符号差}
为 $(2, 1)$，与第 17 题一致，无需\omterm{def:b2:reduction:eigen}{特征值}
计算。高斯配方法丢失度量数据：新坐标
不是标准正交的，因此\omterm{def:b2:reduction:eigen}{特征值}（曲面的形状、
其轴及轴长）均已丢失；只剩
仿射类型。

\textbf{24.} 设 $p$、$n$、$z$ 分别为正、
负与零\omterm{def:b2:reduction:eigen}{特征值}的个数，$p + n + z = 3$。笛卡儿
符号法则给出 $p$ 的上界 $V$，即
$\chi_A$ 的变号次数；同理，$n$ 的上界为 $V'$，即
$\chi_A(-\lambda)$ 的变号次数；此外，每对相邻非零系数
恰在两个多项式之一中产生一次变号，
因而 $V + V' \leq 3 - z$（零根表现为末尾
的零系数）。于是 $p + n = 3 - z \geq
V + V' \geq p + n$，等号成立，故 $p = V$ 恰成立——即
\omterm{def:b2:reduction:eigen}{特征值}的符号可由此读出。对 $\chi_A(\lambda)
= \lambda^3 - 3\lambda^2 - 9\lambda - 5$：符号为 $+,-,-,-$，
得 $V = 1$，而 $\chi_A(-\lambda) = -\lambda^3 - 3\lambda^2
+ 9\lambda - 5$ 的符号为 $-,-,+,-$，$V' = 2$。\omterm{thm:b2:quadratic:sylvester}{符号差} $(1,
2)$，与第 6 题的精确分解 $\chi_A =
(\lambda - 5)(\lambda + 1)^2$ 一致。

\textbf{25.} 算法：对二次部分作标准正交对角化
（谱定理：这是唯一在分析上较深的一步，
也正是它使分类具有
\emph{欧氏}性质）；求解 $A\Omega = -b$ 以
判定中心型与抛物型，并在可能时
消去线性部分（仿射几何）；由
秩、\omterm{thm:b2:quadratic:sylvester}{符号差}与常数项读出类型（西尔维斯特
惯性定理保证这些量为不变量）。若只关心仿射类型，
则高斯配方法可代替谱定理，
代价是失去度量信息。在 $\R^2$ 中，同一机制
即可对二次曲线分类：椭圆、双曲线、抛物线，
以及一对直线、一条直线、一个点与空集。在 $\R^n$
中的变化只在于记账：类型由
$A$ 的\omterm{thm:b2:quadratic:sylvester}{符号差}、$b$ 相对于
$\operatorname{im}A$ 的位置及一个常数决定，其中
$(n{+}2)$ 维的 $(A, b, c)$ 加边矩阵
给出一个紧凑的不变量。
\end{solution}
